\documentclass[11pt,a4paper]{article}

% ---------------------------------------------------------
% PACKAGES
% ---------------------------------------------------------
\usepackage[utf8]{inputenc}
\usepackage[T1]{fontenc}
\usepackage{lmodern}
\usepackage{amsmath, amssymb, amsfonts}
\usepackage{bm}
\usepackage{graphicx}
\usepackage{hyperref}
\usepackage{authblk}
\usepackage{geometry}
\usepackage{cite}
\usepackage{physics}
\usepackage{tensor}
\usepackage{enumitem}

\geometry{margin=1in}

% ---------------------------------------------------------
% TITLE
% ---------------------------------------------------------
\title{\textbf{Companion CI V3.0: Black Holes and White Holes as CPT-Conjugate\\
Inter-Sector Horizons in the Relaxation-Driven Cyclic Cosmology (RDCC)}}

\author[1]{Michael Lehmann}
\affil[1]{\small Independent Researcher, Relaxation-Driven Cyclic Cosmology (RDCC) Framework\\
\texttt{mi.lehmann@gmx.de}}

\date{\small \today}

% ---------------------------------------------------------
% DOCUMENT START
% ---------------------------------------------------------
\begin{document}

\maketitle

% ========================================================
% ABSTRACT
% ========================================================

\begin{abstract}
In the maximally extended Schwarzschild geometry a white hole is the time-reversed
counterpart of a black hole. In the Relaxation-Driven Cyclic Cosmology (RDCC), this
mathematical duality acquires a physical interpretation. Because the global cosmological
state resides in a bipartite Hilbert space and is exactly invariant under a global CPT
transformation, the visible sector $(+)$ and the mirror sector $(-)$ are CPT conjugates.
A black hole in the $(+)$ sector corresponds to a white hole in the $(-)$ sector, and the
apparent one-way nature of horizons is a sectoral projection of a globally symmetric
inter-sector interface.

The minimal dimension-six inter-sector operator and the relaxation dynamics governed
by the parameter $\alpha$ regulate horizon microphysics, smooth classical singularities,
and ensure global information conservation. Black holes and white holes emerge as two
sectoral projections of a single CPT-symmetric inter-sector horizon. This Companion
develops the unified RDCC interpretation of horizon physics and outlines observational
implications for gravitational-wave echoes, the stochastic gravitational-wave background,
and weak lensing around supermassive black holes.
\end{abstract}

% ========================================================
% 1 INTRODUCTION
% ========================================================

\section{Introduction}

Black holes are among the most robust predictions of General Relativity, supported by
observations of X-ray binaries, gravitational-wave mergers, and horizon-scale imaging.
White holes, which appear in the maximally extended Schwarzschild solution as the
time-reversed counterparts of black holes, are usually regarded as unphysical artifacts of
analytic continuation.

The Relaxation-Driven Cyclic Cosmology (RDCC) offers a different perspective. RDCC
postulates that the global cosmological state lives in a bipartite Hilbert space,
\begin{equation}
\mathcal{H}_{\mathrm{global}} = \mathcal{H}_{+} \otimes \mathcal{H}_{-},
\end{equation}
and is exactly invariant under a global CPT transformation that exchanges the two
sectors. Observers are confined to one sector and only have access to the reduced density
matrix obtained by tracing out the other sector. Causal structure, entropy flow, and the
arrow of time are therefore sectoral phenomena.

In this setting the black-hole/white-hole duality of the maximally extended Schwarzschild
geometry acquires a natural physical interpretation. A black hole in the $(+)$ sector is
mapped by CPT to a white hole in the $(-)$ sector. The one-way nature of horizons is
not an absolute feature of spacetime but a consequence of sectoral projection. From the
global point of view the horizon is an inter-sector interface across which information is
redistributed but not destroyed.

This Companion develops this interpretation systematically. Section~2 reviews the RDCC
framework and the bipartite Hilbert-space structure. Section~3 analyzes the Schwarzschild
geometry and its CPT mapping. Section~4 develops the sectoral-projection interpretation
of horizon causality. Section~5 introduces the inter-sector operator and relaxation
dynamics near horizons. Section~6 explains global information conservation and the
resolution of the black-hole information paradox. Section~7 outlines observational
implications, and Section~8 connects horizon physics to the thermodynamic arrow of
time and the cyclic structure of RDCC.

% ========================================================
% 2 RDCC FRAMEWORK AND BIPARTITE HILBERT SPACE
% ========================================================

\section{RDCC Framework and Bipartite Hilbert Space}

The Relaxation-Driven Cyclic Cosmology (RDCC) is built on a single structural postulate:
the complete cosmological quantum state resides in a bipartite Hilbert space,
\begin{equation}
\mathcal{H}_{\mathrm{global}} = \mathcal{H}_{+} \otimes \mathcal{H}_{-},
\end{equation}
and is exactly invariant under a global CPT transformation. The visible sector $(+)$ and
the mirror sector $(-)$ are exchanged by CPT,
\begin{equation}
\ket{\Psi_{\mathrm{global}}}
=
\mathrm{CPT}\,\ket{\Psi_{\mathrm{global}}},
\qquad
\mathrm{CPT}:\mathcal{H}_{+}\rightarrow\mathcal{H}_{-}.
\end{equation}

Observers are confined to one sector and only have access to the reduced density matrix
obtained by tracing out the other sector. For an observer in the $(+)$ sector the relevant
state is
\begin{equation}
\rho_{+} = \mathrm{Tr}_{-}\!\left[\,\ket{\Psi_{\mathrm{global}}}\bra{\Psi_{\mathrm{global}}}\,\right].
\end{equation}
Sectoral projection breaks the manifest CPT symmetry of the global state and produces
an apparent arrow of time, entropy increase, and irreversible macroscopic evolution.

This structure applies not only to cosmological degrees of freedom but also to local
gravitational configurations such as black holes. A black hole in the $(+)$ sector is part of
the global state and has a CPT-conjugate counterpart in the $(-)$ sector. The horizon is
therefore naturally interpreted as an inter-sector interface rather than a purely intra-sector
boundary.

RDCC introduces a minimal inter-sector coupling through a dimension-six operator of the
schematic form
\begin{equation}
\mathcal{O}_{\mathrm{int}}
\;\sim\;
\frac{1}{M^{2}}
\,T^{(+)}_{\mu\nu}\,T^{(-)\,\mu\nu},
\end{equation}
where $T^{(\pm)}_{\mu\nu}$ are the stress-energy tensors in the two sectors and $M$ is a
high-energy scale. This operator mediates weak but nonzero correlations between the
sectors. Its effective impact is controlled by the relaxation parameter $\alpha$, which also
governs the cyclic cosmological dynamics.

In the cosmological context this structure leads to a nonsingular bounce, a natural
resolution of the Tolman low-entropy problem, and correlated small deviations from the
standard cosmological model. In the present Companion we apply the same structural
ingredients to black-hole spacetimes and show that the horizon becomes a CPT-symmetric
inter-sector interface across which information is redistributed but not destroyed.

% ========================================================
% 3 SCHWARZSCHILD GEOMETRY AND CPT TIME REVERSAL
% ========================================================

\section{Schwarzschild Geometry and CPT Time Reversal}

The exterior Schwarzschild metric in standard coordinates is
\begin{equation}
ds^{2}
=
-\left(1-\frac{2M}{r}\right) dt^{2}
+
\left(1-\frac{2M}{r}\right)^{-1} dr^{2}
+
r^{2} d\Omega^{2},
\end{equation}
where $M$ is the mass parameter and $d\Omega^{2}$ is the metric on the unit two-sphere.
This coordinate patch covers the region outside the event horizon at $r = 2M$.

To reveal the full causal structure one introduces Kruskal–Szekeres coordinates $(U,V)$,
in which the maximally extended Schwarzschild solution becomes analytic and contains
four regions: the exterior region I, the black-hole interior II, the white-hole interior III,
and a second exterior region IV. The event horizons are null surfaces separating these
regions.

Time reversal exchanges the black-hole and white-hole regions. A black hole is
characterized by a future event horizon from which no signal can escape, while a white
hole has a past event horizon into which no signal can enter. Both structures are allowed
by the classical field equations, but only black holes are considered physically realized in
standard treatments.

In the RDCC framework this asymmetry is not fundamental. The global CPT symmetry
of the cosmological state implies that for every black-hole configuration in the $(+)$ sector
there exists a CPT-conjugate configuration in the $(-)$ sector. In the maximally extended
Schwarzschild geometry this conjugate configuration is naturally identified with the
white-hole region.

The full Kruskal diagram can therefore be reinterpreted as a bipartite structure: one
exterior region and its associated black-hole interior belong to the $(+)$ sector, while the
other exterior region and its associated white-hole interior belong to the $(-)$ sector. The
two horizons are sectoral projections of a single CPT-symmetric inter-sector interface.

This reinterpretation preserves the classical geometry while embedding it into the
bipartite Hilbert-space structure of RDCC. The black-hole/white-hole duality becomes a
manifestation of global CPT symmetry rather than a mathematical artifact of analytic
continuation.

% ========================================================
% 4 SECTORAL PROJECTION AND RELATIVE CAUSALITY
% ========================================================

\section{Sectoral Projection and Relative Causality}

In RDCC observers are confined to one sector and only have access to the sectoral
projection of the global state. This has profound consequences for the interpretation of
causal structure near horizons.

Consider a global configuration in which the horizon is a CPT-symmetric inter-sector
interface. From the perspective of the global state, information flows across this interface
in a way that preserves unitarity and CPT symmetry. However, an observer in the $(+)$
sector only sees the projection of this process onto $\mathcal{H}_{+}$. For such an observer
the horizon appears as a one-way surface: nothing can escape from the black-hole interior
back to the exterior region.

From the perspective of the $(-)$ sector the same global process appears time-reversed.
The object corresponding to the black hole in the $(+)$ sector is seen as a white hole in
the $(-)$ sector. For an observer in the $(-)$ sector nothing can enter this object from the
exterior, but signals can emerge from it. The causal structure is therefore relative to the
sectoral projection.

This leads to the following picture:

\begin{itemize}
\item In the $(+)$ sector the horizon is a future event horizon of a black hole.
\item In the $(-)$ sector the same global interface appears as a past event horizon of a white hole.
\end{itemize}

The one-way nature of horizons is not an absolute property of spacetime but a consequence
of restricting attention to a single sector. The global state sees a symmetric interface,
while each sector sees an asymmetric horizon.

This interpretation is fully consistent with the RDCC view of the thermodynamic arrow
of time. Entropy increase and irreversibility are sectoral phenomena that arise from
projection and relaxation. The same mechanism that produces a time-oriented cosmology
in each sector also produces a time-oriented horizon structure.

In this sense, black-hole and white-hole horizons are complementary sectoral projections
of a single CPT-symmetric inter-sector interface. The apparent causal asymmetry is a
projection effect rather than a fundamental feature of the underlying geometry.

% ========================================================
% 5 INTER-SECTOR OPERATOR AND RELAXATION AT HORIZONS
% ========================================================

\section{Inter-Sector Operator and Relaxation at Horizons}

The minimal inter-sector operator of RDCC,
\begin{equation}
\mathcal{O}_{\mathrm{int}}
\;\sim\;
\frac{1}{M^{2}}
\,T^{(+)}_{\mu\nu}\,T^{(-)\,\mu\nu},
\end{equation}
couples the stress-energy content of the two sectors. Far from strong-gravity regions this
coupling is extremely weak, but near black-hole horizons the local energy densities and
curvature invariants become large, enhancing the effective impact of the operator. The
relaxation dynamics governed by the parameter $\alpha$ then play a crucial role.

We can distinguish three dynamical regimes:

\begin{enumerate}
\item \textbf{Far from the horizon:}  
The inter-sector coupling is negligible, and the spacetime is well described by the
classical Schwarzschild solution in each sector. Local observables are effectively
decoupled between sectors.

\item \textbf{Near the horizon:}  
Strong curvature amplifies the effective inter-sector coupling. The relaxation dynamics
drive the local state toward a configuration in which the classical singularity is replaced
by a smooth, nonsingular region. The horizon becomes an inter-sector interface across
which information is redistributed.

\item \textbf{Deep in the interior:}  
The classical singularity is avoided. Instead, the global state undergoes a bounce-like
transition in which degrees of freedom are reconfigured between the sectors. The
relaxation timescale
\begin{equation}
\tau_{\mathrm{rel}} \sim \frac{1}{\alpha}
\end{equation}
controls how long the black hole appears quasi-eternal from the perspective of a
sectoral observer.
\end{enumerate}

The key point is that the inter-sector operator does not introduce new fields or ad-hoc
boundary conditions. It is the same structural ingredient that appears in the cosmological
context. The role of $\alpha$ is also unified: it governs both the large-scale relaxation that
drives the cyclic cosmology and the local relaxation that smooths out horizon
singularities.

% --------------------------------------------------------
% 5.1 No Classical Matter Transmission Across the Interface
% --------------------------------------------------------

\subsection{No Classical Matter Transmission Across the Interface}

The inter-sector operator mediates correlations between the $(+)$ and $(-)$ sectors, but
it does not permit coherent transmission of classical matter across the inter-sector
horizon. The operator couples stress-energy tensors and therefore acts on the quantum
state rather than on individual macroscopic configurations.

Near the horizon the local state becomes highly mixed due to strong curvature and
relaxation effects. Any infalling macroscopic object is effectively decohered and
thermalized before its information is redistributed globally. The global CPT-symmetric
evolution therefore transfers information, not matter, across the interface.

This distinction is essential: the inter-sector horizon is not a traversable geometric
bridge. It is a locus where the global bipartite state reorganizes its degrees of freedom.
Classical matter does not emerge in the $(-)$ sector as a CPT-transformed counterpart,
and no dynamical annihilation processes are triggered. Only the informational content of
the collapsing state is encoded in the mirror-sector projection.

% ========================================================
% 6 GLOBAL INFORMATION CONSERVATION AND THE BLACK-HOLE INFORMATION PARADOX
% ========================================================

\section{Global Information Conservation and the Black-Hole Information Paradox}

In semiclassical gravity the black-hole information paradox arises because Hawking
radiation appears thermal and uncorrelated with the initial state. If the black hole
evaporates completely, the final state seems mixed even if the initial state was pure,
apparently violating unitarity.

In RDCC the fundamental object is the global state $\ket{\Psi_{\mathrm{global}}}$ in the
bipartite Hilbert space. This state evolves unitarily under a CPT-symmetric Hamiltonian
that includes the inter-sector operator. Information is never lost at the global level.
What appears as information loss in one sector is compensated by information flow into
the other sector.

From the perspective of an observer in the $(+)$ sector the formation and evaporation of
a black hole resemble the standard picture: matter collapses, a horizon forms, and
Hawking-like radiation is emitted. However, the detailed structure of the radiation is
modified by the inter-sector coupling. Correlations that would be invisible in a purely
intra-sector description are encoded in entanglement patterns between the emitted
quanta and degrees of freedom in the $(-)$ sector.

The apparent loss of information in the $(+)$ sector is therefore an artifact of tracing out
the $(-)$ sector. If one had access to the full global state, one would see that the process
is unitary. The black-hole information paradox is thus reinterpreted as a \emph{projection
paradox}: it arises from insisting on a single-sector description of a fundamentally
bipartite process.

This perspective also clarifies the role of white holes. In the global picture the black-hole
horizon in the $(+)$ sector and the white-hole horizon in the $(-)$ sector are two aspects
of the same inter-sector interface. Information that appears to fall irreversibly into the
black hole from the $(+)$ perspective can re-emerge in the $(-)$ sector as white-hole-like
emission. The global CPT symmetry ensures that the full process is time-symmetric and
unitary, even though each sector sees a time-oriented, apparently irreversible evolution.

A further structural consequence of the bipartite Hilbert-space formulation is that a black
hole in the $(+)$ sector and its CPT-conjugate white hole in the $(-)$ sector do not
represent two independent objects. They form a single global configuration whose
apparent separation arises only from sectoral projection. From the viewpoint of the
global state the horizon is therefore not a terminal boundary but an inter-sector interface
across which information is redistributed in a CPT-symmetric manner.

This implies that the black-hole/white-hole pair acts as a global information bridge
between the two sectors. Energy and information are conserved globally, but each sector
sees only its own projection of the process. What appears as irreversible infall in the $(+)$
sector corresponds to time-reversed outflow in the $(-)$ sector. The two processes are
complementary aspects of a single CPT-symmetric evolution.

% --------------------------------------------------------
% 6.1 Information vs. Matter Transfer Across Inter-Sector Horizons
% --------------------------------------------------------

\subsection{Information vs. Matter Transfer Across Inter-Sector Horizons}

A common misconception arises when interpreting the CPT-conjugate relation between a
black hole in the $(+)$ sector and its white-hole counterpart in the $(-)$ sector. One might
naively expect that an infalling macroscopic object could ``reappear'' in the mirror sector
as an outgoing object. This interpretation is not supported by the RDCC framework.

The inter-sector horizon does not transmit classical matter or macroscopic structures.
What is transferred across the interface is information encoded in the global bipartite
quantum state. The relaxation dynamics and the dimension-six operator mediate
correlations between sectors, but they do not permit coherent transport of extended
objects. Near the horizon the local state becomes highly mixed, and macroscopic
configurations are effectively thermalized before their information is redistributed
globally.

CPT symmetry acts on the global state, not on individual objects undergoing a dynamical
transition between sectors. Consequently, no object emerges in the $(-)$ sector as a
classical CPT-transformed counterpart. The white-hole emission observed in the $(-)$
sector is a sectoral projection of the global information flow, not a material reappearance
of infalling matter.

This resolves potential inconsistencies such as immediate annihilation of transferred
matter with ambient mirror-sector matter. Since no classical matter is transmitted, only
information, the global CPT-symmetric evolution remains well defined without requiring
direct dynamical transport of macroscopic objects between sectors.

% ========================================================
% 7 OBSERVATIONAL IMPLICATIONS
% ========================================================

\section{Observational Implications}

Although the inter-sector coupling is weak, the combination of strong gravity and long
timescales near black holes opens the possibility of observable signatures. These signatures
arise from modifications of horizon microphysics, relaxation dynamics, and the global
bipartite structure of RDCC. We outline three classes of potential effects.

% --------------------------------------------------------
% 7.1 Gravitational-Wave Echoes
% --------------------------------------------------------

\subsection{Gravitational-Wave Echoes}

If the horizon is not a perfectly absorbing surface but an inter-sector interface with
nontrivial microphysics, the ringdown signal of a perturbed black hole may exhibit
late-time echoes. In RDCC the relaxation dynamics near the horizon introduce a
characteristic timescale,
\begin{equation}
\tau_{\mathrm{rel}} \sim \frac{1}{\alpha},
\end{equation}
which can manifest as a delay between the primary ringdown and subsequent echo-like
features.

The detailed waveform depends on the effective reflectivity of the inter-sector interface
and the frequency dependence of the inter-sector coupling. Qualitatively, one expects
small-amplitude, phase-shifted repetitions of the ringdown signal at intervals set by the
light-travel time between the effective interface and the photon sphere, modulated by the
relaxation dynamics. Future observations with LIGO, Virgo, KAGRA, LISA, and the
Einstein Telescope could constrain such scenarios.

% --------------------------------------------------------
% 7.2 Stochastic Gravitational-Wave Background
% --------------------------------------------------------

\subsection{Stochastic Gravitational-Wave Background}

The cumulative effect of many black-hole formation and merger events in the presence of
inter-sector coupling can imprint characteristic features on the stochastic gravitational-wave
background (SGWB). If horizon microphysics in RDCC modifies the late-time tail of
ringdown signals or introduces additional channels of energy transfer between sectors, the
resulting background may deviate from the standard prediction.

In particular, the same relaxation parameter $\alpha$ that governs cosmological deviations
from the standard model can induce a mild scale dependence in the gravitational-wave
spectrum. This connects the black-hole sector to the cosmological sector and provides a
cross-check of the RDCC framework.

% --------------------------------------------------------
% 7.3 Weak Lensing Around Supermassive Black Holes
% --------------------------------------------------------

\subsection{Weak Lensing Around Supermassive Black Holes}

Inter-sector leakage near horizons can in principle alter the effective stress-energy
distribution in the vicinity of supermassive black holes. While the effect is expected to be
small, it could lead to subtle deviations in the lensing patterns of background sources.
High-resolution imaging and precise weak-lensing measurements around galactic centers
may reveal small anomalies that are difficult to explain within a purely classical framework.

Such signatures would be challenging to isolate, but they illustrate that RDCC does not
confine its predictions to the early universe. The same structural postulate that shapes the
global cosmology also has local consequences in strong-gravity regimes.

% ========================================================
% 8 RELATION TO THE THERMODYNAMIC ARROW OF TIME AND CYCLIC COSMOLOGY
% ========================================================

\section{Relation to the Thermodynamic Arrow of Time and Cyclic Cosmology}

One of the central achievements of the Relaxation-Driven Cyclic Cosmology (RDCC) is the
unification of the thermodynamic arrow of time, the resolution of cosmological singularities,
and the cyclic structure of the universe under a single structural postulate: global CPT
symmetry of the bipartite cosmological quantum state combined with relaxation dynamics.

The present Companion shows that the same ingredients naturally extend to black-hole
physics. The apparent irreversibility of black-hole formation and evaporation is a sectoral
phenomenon, just like the cosmological arrow of time. The classical singularity inside a
black hole is replaced by a nonsingular region in which the global state undergoes a
bounce-like reconfiguration between sectors, analogous to the cosmological bounce.

In this sense black holes are local echoes of the global cyclic structure. They are regions
where the interplay between CPT symmetry, bipartite structure, and relaxation dynamics
becomes manifest on astrophysical scales. The black-hole information paradox is not an
isolated puzzle but part of a broader pattern: whenever one insists on a single-sector
description of a fundamentally bipartite, CPT-symmetric process, apparent paradoxes
arise that dissolve once the global structure is taken into account.

The inter-sector horizon therefore plays a dual role. Cosmologically, it ensures that the
universe evolves through nonsingular cycles driven by relaxation. Locally, it ensures that
black holes evolve through nonsingular interior transitions governed by the same
relaxation parameter $\alpha$. In both cases the arrow of time emerges from sectoral
projection rather than from fundamental microscopic irreversibility.

This unified perspective highlights the conceptual coherence of RDCC: the same structural
postulate that governs the largest scales of the universe also governs the strongest-gravity
regimes. Black holes and cosmological bounces are two manifestations of the same
CPT-symmetric, relaxation-driven dynamics.

% ========================================================
% 9 CONCLUSION
% ========================================================

\section{Conclusion}

In this Companion we have developed a unified RDCC-based interpretation of black holes
and white holes as CPT-conjugate inter-sector horizons. Starting from the maximally
extended Schwarzschild geometry, we showed that the mathematical duality between
black-hole and white-hole regions acquires a physical meaning once the geometry is
embedded into the bipartite, CPT-symmetric Hilbert-space structure of the
Relaxation-Driven Cyclic Cosmology (RDCC).

A black hole in the visible $(+)$ sector corresponds to a white hole in the mirror $(-)$
sector, and the apparent one-way nature of horizons is a projection effect of sectoral
observation. From the global perspective the horizon is not a terminal boundary but a
CPT-symmetric inter-sector interface across which information is redistributed while
global unitarity is preserved.

The minimal dimension-six inter-sector operator and the relaxation dynamics governed by
the parameter $\alpha$ play a central role. Near horizons they induce entanglement
leakage between sectors, smooth out classical singularities, and ensure that the global
state undergoes a nonsingular, bounce-like reconfiguration rather than forming a
singularity. The black-hole information paradox is therefore reinterpreted as a
\emph{projection paradox}: information is never lost globally, but sectoral observers who
trace out the other sector perceive an apparent loss.

We have outlined potential observational signatures, including gravitational-wave echoes,
modifications to the stochastic gravitational-wave background, and subtle weak-lensing
effects around supermassive black holes. While these signatures are challenging to detect,
they illustrate that RDCC connects cosmological and astrophysical phenomena within a
single structural framework.

This Companion extends the reach of RDCC from the largest scales of the universe down
to the strong-gravity regime of black holes. It reinforces the central message of the
program: many long-standing puzzles in cosmology and gravity can be addressed
coherently by combining global CPT symmetry, a bipartite Hilbert space, and relaxation
dynamics governed by a single free parameter.

% ========================================================
% APPENDIX A — TECHNICAL STRUCTURE OF THE INTER-SECTOR OPERATOR
% ========================================================

\appendix

\section*{Appendix A: Technical Structure of the Inter-Sector Operator}

The minimal inter-sector operator used throughout RDCC has the schematic form
\begin{equation}
\mathcal{O}_{\mathrm{int}}
=
\frac{1}{M^{2}}
\,T^{(+)}_{\mu\nu}\,T^{(-)\,\mu\nu},
\end{equation}
where $T^{(\pm)}_{\mu\nu}$ denote the stress-energy tensors in the visible and mirror
sectors. This operator is the unique dimension-six scalar that couples the two sectors
without introducing new fields or violating CPT symmetry.

The operator acts on the global bipartite state,
\begin{equation}
\mathcal{H}_{\mathrm{global}}
=
\mathcal{H}_{+}\otimes\mathcal{H}_{-},
\end{equation}
and induces weak but nonzero correlations between the sectors. Its effective strength is
controlled by the relaxation parameter $\alpha$, which governs both cosmological and
local strong-gravity dynamics.

Near black-hole horizons the local curvature enhances the effective contribution of
$\mathcal{O}_{\mathrm{int}}$, leading to increased inter-sector entanglement and smoothing of
the classical singularity. Far from strong-gravity regions the operator is negligible, and
the sectors evolve effectively independently.

The operator does not mediate classical matter transfer. It acts on the quantum state,
redistributing information across the inter-sector interface while preserving global
unitarity and CPT symmetry.


\end{document}
